\chapter{Package Token Policy and Reflection}
\label{ch:v151-package-token-policy}

The v1.5 layer replaced the function-like token notation \(\Tok_{\Pi}(s)\) by the judgment \(\TokIntro(\Pi,s,p)\). This chapter packages the necessary token behavior into a single proof interface.

\begin{definition}[Package token policy]
For a probe bundle \(\Pi\), a package token policy, written
\[
  \PkgTokPol(\Pi),
\]
consists of the following clauses for the concrete signature-token instance.
\begin{enumerate}
  \item \textbf{Token existence.} Every generated answer history \(s\) that is intended to be visible admits some token \(p\) with \(\TokIntro(\Pi,s,p)\).
  \item \textbf{Soundness.} If \(\TokIntro(\Pi,s,p)\), \(\TokIntro(\Pi,t,q)\), and \(\hsame(s,t)\), then \(\psame(p,q)\).
  \item \textbf{Reflection.} If \(\TokIntro(\Pi,s,p)\), \(\TokIntro(\Pi,t,q)\), and \(\psame(p,q)\), then \(\hsame(s,t)\).
  \item \textbf{Package-sameness closure.} On introduced tokens, \(\psame\) is closed under reflexivity, symmetry, and transitivity.
  \item \textbf{No external token equality.} Visible tokens are compared by \(\psame\), not by any imported equality judgment.
\end{enumerate}
\end{definition}

\begin{theorem}[Package-sameness is equivalent to signature-sameness]
\label{thm:v151-psame-iff-hsame}
Assume \(\PkgTokPol(\Pi)\). If
\[
  \TokIntro(\Pi,s,p),
  \qquad
  \TokIntro(\Pi,t,q),
\]
then
\[
  \psame(p,q)
  \quad\Longleftrightarrow\quad
  \hsame(s,t).
\]
\end{theorem}

\begin{proof}
The implication from right to left is the soundness clause of \(\PkgTokPol(\Pi)\). The implication from left to right is the reflection clause. No external token equality is used.
\end{proof}

\begin{corollary}[Concrete package completeness is policy-grounded]
The completeness direction of the concrete signature Globalize instance is frozen only under \(\PkgTokPol(\Pi)\). In particular, it is not a theorem for arbitrary package-sameness relations.
\end{corollary}

\begin{proof}
Completeness requires recovering \(\hsame(s,t)\) from \(\psame(p,q)\). This is exactly the reflection field of \(\PkgTokPol(\Pi)\).
\end{proof}

\begin{remark}[Proof-obligation boundary]
The token policy is deliberately narrow. It applies to the concrete signature-token package instance. Other package constructions must provide their own soundness and reflection certificates before they may inherit the same completeness theorem.
\end{remark}
